
\section{The VALOR SDK}

The development of a {\em physics parameterization} is required
for the analysis and interpretation of any physics dataset.
Typically, this parameterization takes the form of a likelihood function, $\lambda$,
which depends on some aspects of the recorded data and, 
in general, is a function of 
M physics parameters $\vec{\theta}$ 
(parameters of the physics hypothesis used for interpreting the data)
and
N nuisance (systematic) parameters $\vec{f}$.
A successful generic analysis framework should facilitate the implementation of the analysis-specific choices on which aspects of a recorded dataset are to be included in the analysis (such as, for example, which kinematic distributions from which topological selections applied on which detector and/or running conditions), and it should allow the construction of corresponding predictions and the efficient parameterization of the effects of
different physics hypotheses as of systematic effects on these predictions. 
In general, these analysis-specific choices include:
\begin{itemize}
    \item The definition of topological event samples s.
          In principle, different topological samples can be used for 
          different detectors included in the analysis,
          allowing, for example, the higher statistics near detector samples to be
          further subdivided into several exclusive channels and
          utilise correlations in order to disentangle systematic effects.
    \item The definition of reconstructed and true kinematic spaces.
          This includes the choice of kinematic space dimensionality,
          the choice of kinematic quantities, and the choice of binning schemes. 
          In principle, the choice of a reconstructed kinematic space
          can be unique to each topological sample included in the fit.
    \item The choice of generator-level labels (modes m), whose contributions to
          the predicted event rate for a given topological sample are individually 
          tracked. This choice can be unique to each topological sample included 
          in the fit and it is crucial (along with the choice of kinematic quantities)
          so that the application of systematic and physics effects can be 
          finally targeted and expressed with the most natural degrees of freedom.
    \item The choice of systematic parameters, the parameterization of the 
          effect of systematic parameters on the computed event rates, and the
          definition of prior systematic parameter constraints.
\end{itemize}

The general framework used in the VALOR SDK for constructing that physics parameterization, is outlined in Sec. \ref{S:valor_physics_parameterization}.

\subsection{Framework for construction of the likelihood}
\label{S:valor_physics_parameterization}

In the VALOR SDK an arbitrary number of samples, each corresponding
to a given i) detector (or sub-detector region), ii) beam configuration,
iii) observed final state, and iv) kinematic phase space,
can be fit jointly to determine the parameters of a physics hypothesis
in the presence of uncertainties.

A crucial element of parameter and error estimation in the context of
an oscillation analysis
is the efficient calculation of event rate predictions,
both nominal and varied ones corresponding to specific combination of
systematic effects and oscillation (or other new physics) hypotheses.
Event rate predictions needs to be calculated for each topological event
sample included in the fit, as a function of the observed kinematic
quantities chosen for each sample.
Predictions are constructed from Monte Carlo (MC) templates T,
which represent the number of MC events 
(after event selection cuts were applied on the output 
of the full event simulation and reconstruction chain) 
as function of several reconstructed and true quantities,
as needed in order to apply physics and systematic effects
as function of the most appropriate variables, 
and in order to enable comparisons with the chosen
observed kinematic distributions.
Separate MC templates are constructed for each detector d, beam configuration b, 
and sample s. Each MC template contains information on how the number of
events is distributed in the same $N_{r}$-dimensional space {\bf $K_{r}$}  
of reconstructed kinematic variables chosen for the fit samples.
The same reconstructed kinematic variable binning scheme is used
both for the MC templates and the fit distributions.
The MC templates provide a mapping between reconstructed and true information.
Separate templates are constructed for different true reaction modes m, and 
each template contains information on how the number of events, 
for each individual reconstructed bin,
is distributed in a chosen $N_{t}$-dimensional space {\bf $K_{t}$}  
of true kinematic variables.
The true reaction modes, the true kinematic variables, and the kinematic variable
binning schemes are defined so that the intended flavour dependencies, 
reaction mode dependencies and  kinematic dependencies of systematics 
and/or of considered physics hypotheses can be taken into account. 
Therefore, summarizing all MC template dependencies, we can write
\begin{equation}
   T = T_{d;b;s;m}(r,t)
\label{eq:template_deps}
\end{equation}
where r is a bin in the {\bf $K_{r}$} space of reconstructed kinematic variables, and
t is a bin in the {\bf $K_{t}$}  space of true kinematic variables.
A MC template $T_{d;b;s;m}$ is constructed from a MC sample corresponding to an
exposure of $e_{d;b}^{MC}$ and then used to predict observations, or fit data,
corresponding to an exposure of $e_{d;b}^{data}$. Here, we will assume that a
scaling factor $e_{d;b}^{data}$ / $e_{d;b}^{MC}$ is absorbed in the definition of $T_{d;b;s;m}$.

Using MC templates, the predicted number of events 
$n_{d;b;s}^{pred}(r; \vec{\theta}; \vec{f})$ 
in a $N_{r}$-dimensional reconstructed kinematic bin r, 
for a specific sample s, seen in a detector d exposed in a beam configuration b,
and for a particular set of 
M physics parameters $\vec{\theta}$ = $(\theta_0, \theta_1,..., \theta_{M-1})$ and
N nuisance (systematic) parameters $\vec{f}$ = $(f_0, f_1, ..., f_{N-1})$,
is computed as
\begin{equation}
n_{d;b;s}^{pred}(r; \vec{\theta}; \vec{f}) =
   \sum_{m} \sum_{t}  P_{d;b;m}(t; \vec{\theta}) \cdot R_{d;b;s;m}(r,t; \vec{f}) \cdot T_{d;b;s;m}(r,t)
\label{eq:valor_npred}
\end{equation}
where 
$P_{d;b;m}(t; \vec{\theta})$ encapsulates the effect of a physics hypothesis (e.g. neutrino oscillations in a 3-flavour framework), and $R_{d;b;s;m}(r,t; \vec{f})$ 
parameterizes the response of a template bin to systematic variations.
In principle, the range of m values in the above sum depends on the sample s.
In addition, the number, type and dimensionality of true bins t is a function of both s and m. 
The above will be assumed implicitly and not written explicitly.

Once the construction of VALOR event rate predictions $n_{d;b;s}^{pred}(r; \vec{\theta}; \vec{f})$ is implemented, through the factorization described in 
Eq. \ref{eq:valor_npred}, physics is extracted through a comparison with 
observed (or simulated fake) data, which are denoted as $n_{d;b;s}^{obs}(r)$
and represent the observed event rate
for each detector d, beam configuration b, topological sample s
and (multi-dimensional) reconstructed kinematic bin r.
VALOR, typically, uses a binned likelihood ratio method and, therefore,
attempts to minimize the quantity:
\begin{equation}
 -2 ln\lambda(\vec{\theta}; \vec{f}) = 
    -2 ln\lambda_{0}(\vec{\theta}; \vec{f}) 
    -2 ln\lambda_{p}(\vec{\theta}; \vec{f})
\label{eq:valor_likelihood_ratio}
\end{equation}
where 
\begin{equation}
 -2 ln\lambda_{o}(\vec{\theta}; \vec{f}) = 2 \sum_{d,b,s,r}
    \Big( 
       n_{d;b;s}^{obs}(r) \cdot
         ln \frac{n_{d;b;s}^{obs}(r)}{n_{d;b;s}^{pred}(r; \vec{\theta}; \vec{f})} 
       + (n_{d;b;s}^{pred}(r; \vec{\theta}; \vec{f}) - n_{d;b;s}^{obs}(r))
    \Big) 
\label{eq:valor_likelihood_ratio_data}
\end{equation}
and $-2 ln\lambda_{p}$ is a penalty term encapsulating prior constraints
from external data and theory.
The advantage of the likelihood ratio method, 
compared with the extended maximum likelihood method, is that in the
large-sample limit, the quantity $-2 ln\lambda(\vec{\theta}; \vec{f})$
has a $\chi^2$ distribution, and therefore it can be used as a goodness of fit test.
Various strategies are employed for the minimisation of 
$-2 ln\lambda(\vec{\theta}; \vec{f})$, the elimination of nuisance and/or
interesting physics parameters, and the construction of confidence intervals.
These statistical procedures implemented within the VALOR SDK are
summarised in Sec. \ref{S:valor_stat}.


\subsection{Computing effects of physics hypotheses}

The effect of a physics hypothesis, such as, for example, neutrino oscillations,
enters via the term $P_{d;b;m}(t; \vec{\theta})$ in Eq. \ref{eq:valor_npred}.
This term is naturally only a function of true 
kinematic variables and of neutrino flavour and/or true reaction mode 
(both of which are encapsulated in m). For multi-detector fits, 
the dependence on d and b reflects the dependence of $P$ on the appropriate 
baseline (or on the distribution of  baselines, where that is appropriate).

The oscillation probability $P_{d;b;m}(t; \vec{\theta})$ is evaluated in one
of a number of different physics frameworks, each with its own set of parameters $\vec{\theta}$. Currently, the VALOR SDK supports oscillation probability calculations 
in a 3-flavour framework, as well as 3+1 and 3+2 frameworks, considering matter-effects 
in constant density matter and, optionally, non-standard interactions. 
In addition, simpler 2-flavour oscillation probability calculations are also supported.
The application of the term $P_{d;b;m}(t; \vec{\theta})$ requires some extra elaboration

%
%
%

\subsubsection{Standard 3-flavour oscillations}
\label{sec:3fosc}

%
%
\subsubsection*{3-flavour neutrino oscillation probabilities in vacuum}
 
The 3-flavour neutrino oscillation probabilities in vacuum are calculated from equation 11 in \cite{kayser2005}:
\begin{equation}
P(\nu_{\alpha} \rightarrow \nu_{\beta}) = \delta_{\alpha \beta} - 4 \sum_{i>j} Re(U_{\alpha i}^{*} U_{\beta i} U_{\alpha j} U_{\beta j}^{*})
\sin^{2} \left( \frac{\Delta m_{ij}^{2} L}{4E} \right) + 2 \sum_{i>j} Im(U_{\alpha i}^{*} U_{\beta i} U_{\alpha j} U_{\beta j}^{*})
\sin \left(\frac{\Delta m_{ij}^{2} L}{2E} \right)
\label{3flavouroscprobvacuum}
\end{equation}
In equation \ref{3flavouroscprobvacuum}, $\Delta m_{ij}^{2}$ = $m_{j}^{2}$ - $m_{i}^{2}$, L is the neutrino baseline, E is the neutrino energy, and U is the PMNS matrix in vacuum:

\begin{equation}
U = \left(
\begin{array}{ccc}
c_{12}c_{13} & c_{13}s_{12} & s_{13}e^{-i\delta} \\
-c_{23}s_{12} - s_{13}s_{23}c_{12}e^{i\delta} & c_{12}c_{23} - s_{12}s_{13}s_{23}e^{i\delta} & c_{13}s_{23} \\
s_{12}s_{23} - s_{13}c_{12}c_{23}e^{i\delta} & -s_{23}c_{12} - s_{12}c_{23}s_{13}e^{i\delta} & c_{13}c_{23}
\end{array} \right)
\label{PMNSmatrix}
\end{equation}
with $s_{12} = \sin(\theta_{12})$, $c_{12} = \cos(\theta_{12})$, etc. $P(\bar{\nu}_{\alpha} \rightarrow \bar{\nu}_{\beta})$
is calculated by changing the sign of the third term in equation \ref{3flavouroscprobvacuum} (this assumes that CPT is
conserved).

%
%
\subsubsection*{3-flavour neutrino oscillation probabilities in constant density matter}

All three neutrino flavours undergo neutral-current interactions with protons, neutrons and electrons in matter. These
neutral-current interactions have identical amplitudes for all three flavours, and they produce no observable effects on
the probabilities of oscillation between one flavour and another. However the oscillation probabilities are changed by
coherent forward scattering of electron neutrinos in charged-current interactions with electrons in matter \cite{kayser2005}.

Using natural units with $\hbar$ = c = 1, the time evolution of neutrino flavour states in vacuum is described by the
Schr\"odinger equation:

\begin{equation}
i \frac{\partial \mathbf{\psi}}{\partial t} = H_{F} \mathbf{\psi}
\end{equation}
where
\begin{equation}
  \mathbf{\psi} =\left(
\begin{array}{c}
\psi_{e} \\
\psi_{\mu} \\
\psi_{\tau}
\end{array} \right)
\end{equation}
and $H_{F}$ is the effective Hamiltonian in the flavour-state basis
\begin{equation}
H_{F} = \frac{1}{2E} U M U^{\dagger}
\end{equation}
with M defined by
\begin{equation}
M = \left(
\begin{array}{ccc}
0 & 0 & 0 \\
0 & \Delta m_{21}^{2} & 0 \\
0 & 0 & \Delta m_{31}^{2}
\end{array} \right)
\end{equation}
The calculation of oscillation probabilities requires neutrino flight time to be converted to distance travelled,
and this conversion assumes that neutrinos are highly relativistic. In
this calculation, $U M U^{\dagger}$ is used for neutrinos, but its complex conjugate is used for antineutrinos. For neutrinos,
the effects of the charged-current interactions between electron neutrinos and electrons are taken into account by adding the
potential $2E \sqrt{2} G_{F} N_{e}$ to the real part of the first diagonal element of $U M U^{\dagger}$, where E is the
neutrino energy, $G_{F}$ is the Fermi coupling constant, and $N_{e}$ is the number density of electrons in matter. For
antineutrinos, this potential is subtracted from the real part of the first diagonal element of the complex conjugate of
$U M U^{\dagger}$ (for further details see section 9.2 in \cite{giunti}).

After the addition or subtraction of this potential, $U M U^{\dagger}$ (or its complex conjugate) is diagonalised. The eigenvalues
of a Hermitian matrix are always real; they are calculated by solving the (cubic) characteristic equation using the method
of del Ferro, Tartaglia and Cardano as described in \cite{kopp}. The differences between the eigenvalues are the effective
mass-squared differences in matter. The eigenvectors are calculated using an algebraic method: one of the components is set
to be 1.0 (real), and the other two components are calculated using two of the three simultaneous equations in
$U M U^{\dagger} - \lambda I = 0$, where the $\lambda$ are the eigenvalues. After being normalised, the three eigenvectors
become the columns of the effective mixing matrix in matter.

The initial flavour state of the neutrino is represented by a 1$\times$3 column vector whose entries are complex; if, for
example, the initial flavour state is a muon neutrino, this column vector is

\begin{displaymath}
\left(
\begin{array}{c}
0 \\
1 \\
0
\end{array} \right)
\end{displaymath}
The vector representing the initial flavour state is then multiplied by the Hermitian conjugate of the matter mixing matrix
to convert it to another 1$\times$3 complex column vector representing the initial mass states. These mass states are propagated
to Super-K by multiplying the jth initial mass state by $\exp \left(\frac{-i \Delta m_{j1}^{2}L}{2E} \right)$, where the
$\Delta m_{j1}^{2}$ are the mass-squared differences in matter, i.e.\ the differences between the eigenvalues of $U M U^{\dagger}$.
This gives a new 1$\times$3 complex column vector representing the final mass states, and these are converted to the final flavour
states by multiplying by the matter mixing matrix. The entries of the resulting 1$\times$3 complex column vector represent the
amplitudes of each flavour at Super-K, and the corresponding oscillation probabilities are calculated as the moduli squared of these
amplitudes.  

%
%
\subsubsection*{Inputs to the 3-flavour neutrino oscillation probability calculations}

When calculating the 3-flavour oscillation probabilities, the mixing angles can be input either in the form $\sin^{2}(2 \theta_{ij})$
or as $\sin^{2}(\theta_{ij})$. $\sin(\theta_{ij})$ and $\cos(\theta_{ij})$ are calculated from the input values, and these are used in
turn to calculate the elements of $U$. The ``solar'' mass-squared difference is input as $\Delta m_{21}^{2}$, and this is $>$ 0 for
both mass hierarchies. The ``atmospheric''
mass-squared difference is input either as $|\Delta m_{32}^{2}|$ or using a definition introduced by G.\ Fogli, E.\ Lisi et al.\ in
\cite{LisiGlobalAnalysis}:
\begin{equation}
\Delta m^{2}_{FL} = m_{3}^{2} - \frac{m_{2}^{2} + m_{1}^{2}}{2}
\end{equation}
The quantity denoted $\Delta m^{2}_{FL}$ is useful since its absolute value is the same for both mass hierarchies, and only its sign
changes. If the atmospheric mass-squared difference is input as $\Delta m^{2}_{32}$, the mass hierarchy is set using the sign of
$\Delta m_{32}^{2}$; this is positive for the normal mass hierarchy (NH) and negative for the inverted mass hierarchy (IH). In this case,
$\Delta m_{31}^{2}$ is not input but is calculated as $\Delta m_{31}^{2}$ = $\Delta m_{21}^{2}$ + $\Delta m_{32}^{2}$ for both hierarchies.
If the atmospheric mass-squared difference is input as $\Delta m^{2}_{FL}$, the mass hierarchy is set using the sign of $\Delta m^{2}_{FL}$.
Then $\Delta m_{32}^{2}$ and $\Delta m_{31}^{2}$ are calculated from the input values of $\Delta m^{2}_{FL}$ and $\Delta m_{21}^{2}$.

%
%
\subsubsection*{Validation of the 3-flavour oscillation probabilities}

The 3-flavour oscillation probabilities in vacuum were checked using an alternative formulation in equations 13.13 and 13.14
in the PDG review \cite{pdgreview}:
\begin{equation}
P(\nu_{\alpha} \rightarrow \nu_{\beta}) = \sum_{j} |U_{\beta j}|^{2} |U_{\alpha j}|^{2} + 2 \sum_{j>k} |U_{\beta j} U_{\alpha j}^{*} U_{\alpha k} U_{\beta k}^{*}|
\cos \left(\frac{\Delta m_{jk}^{2} L}{2E}  - \phi_{\beta \alpha;jk}  \right)
\label{3flavouroscprobvacuumpdg}
\end{equation}
where $\phi_{\beta \alpha;jk}$ = arg($U_{\beta j} U_{\alpha j}^{*} U_{\alpha k} U_{\beta k}^{*}$). $P(\bar{\nu}_{\alpha} \rightarrow \bar{\nu}_{\beta})$
is calculated by changing the sign of $\phi_{\beta \alpha;jk}$ in the argument of the cosine. The probabilities calculated from equation \ref{3flavouroscprobvacuumpdg}
agreed with those calculated from equation \ref{3flavouroscprobvacuum} to 14 significant figures.

Several checks of the oscillation probabilities in constant density matter are written into the code that calculates them:
\begin{enumerate}
\item The eigenvalues of $U M U^{\dagger}$ are checked by comparing their sum with the trace of the matrix.
\item Each normalised eigenvector is multiplied in turn by each of the 3 matrices formed by subtracting the eigenvalues from the real parts of the
diagonal elements of $U M U^{\dagger}$, and a check is made of the differences between these products and a zero vector.
\item A check is made that the normalised eigenvectors of $U M U^{\dagger}$ are orthogonal by calculating the scalar product between them (the scalar product of two
complex vectors is the product of the first vector with the complex conjugate of the second vector).
\item A check is also made that the Hermitian conjugate of the mixing matrix in matter is equal to the inverse of the matrix. This is done by multiplying the
mixing matrix in matter by its Hermitian conjugate, and checking the differences between the product and an identity matrix.
\end{enumerate}

The values of the oscillation probabilities in constant density matter were checked by comparing them with equivalent probabilities calculated by an
independently-written Fortran program \cite{sloan}. This Fortran program uses different algorithms, and calculates numerically the eigenvalues and
eigenvectors of $U M U^{\dagger}$. Nevertheless there was very good agreement between the two calculations of the probabilities, with the
fractional differences being $\approx 2 \times 10^{-6}$.

The values of the matter oscillation probabilities were also compared with those given by Prob3++ \cite{Prob3}, and the fractional differences were
$\approx$ 1-5 $\times$ 10$^{-4}$. The comparison plots can be seen in \cite{Compwithprob3}.

%
%
\subsubsection*{Accuracy of the 3-flavour oscillation probabilities}

As stated in section \ref{3flavourvalidation}, the 3-flavour oscillation probabilities in vacuum agreed to 14 significant figures between the two
formulations.

An estimate of the accuracy of the 3-flavour oscillation probabilities in constant density matter was made in three separate ways:
\begin{enumerate}
\item The matter probabilities for zero density were compared with the vacuum probabilities; there was good agreement, with the fractional differences being
$\approx$ 2-5 $\times$ 10$^{-6}$.
\item The probabilites given when calculating the eigenvectors with simultaneous equations 1 and 2 in $U M U^{\dagger} - \lambda I = 0$ were compared
with the probabilities given when calculating the eigenvectors with simultaneous equations 2 and 3. The fractional differences between these two
calculations were again $\approx$ 2-5 $\times$ 10$^{-6}$.
\item The sum of three matter probabilities, for example P($\nu_{\mu} \rightarrow \nu_{e}$) + P($\nu_{\mu} \rightarrow \nu_{\mu}$) +
P($\nu_{\mu} \rightarrow \nu_{\tau}$), was compared with 1.0; again there was good agreement, and the differences were $\approx 2 \times 10^{-6}$.
\end{enumerate}

This means that the matter oscillation probabilities should be considered to be accurate to 5 decimal places. The accuracy of the matter probabilities
is less than that of the vacuum probabilities due to the method of calculation of the eigenvectors. This involves calculations of the form
\begin{displaymath}
\frac{(ab - cd)}{(ef - gh)}
\end{displaymath}
where a, b, c, etc.\ are elements of $U M U^{\dagger} - \lambda I$. Frequently ab is close in value to cd, which means that (ab - cd) is much smaller than ab or
cd, and the same applies to (ef - gh). However the advantage of the method of calculation of the eigenvectors described in section \ref{matterosc} is that it allows
the worst of these cancellation errors to be avoided; this can be done by calculating the eigenvectors using different pairs of simultaneous equations in
$U M U^{\dagger} - \lambda I = 0$ for different combinations of the oscillation parameters. In practice, the eigenvectors are nearly always calculated
using equations 1 and 2 in $U M U^{\dagger} - \lambda I = 0$, but equations 2 and 3 are used when 3.10 $<$ $\delta_{CP}$ $<$ 3.18 and
$\sin^{2}(2 \theta_{13})$ $>$ $10^{-8}$ in order to avoid such cancellation errors.  

%
%
\subsubsection*{Application of oscillation probabilities to MC templates}

In general, in a 3-flavour analysis,
the CC MC templates constructed from the unoscillated MC samples are 
weighted with the corresponding survival probability: 
The $\nu_{\mu}$ CC templates are weighted with $P(\nu_{\mu}\rightarrow \nu_{\mu})$,
the $\bar{\nu}_{\mu}$ CC templates are weighted with $P(\bar{\nu}_{\mu}\rightarrow \bar{\nu}_{\mu})$,
the $\nu_{e}$ CC templates are weighted with $P(\nu_{e}\rightarrow \nu_{e})$, and
the $\bar{\nu}_{e}$ CC templates are weighted with $P(\bar{\nu}_{e}\rightarrow \bar{\nu}_{e})$,
The CC MC templates made from the {\em oscillated (or swapped)} $\nu_{e}$ event sample
(a $\nu_{e}$ event sample that was generated with the $\nu_{\mu}$ flux, 
as if all $\nu_{\mu}$'s converted to $\nu_{e}$) are weighted with $P(\nu_{\mu}\rightarrow \nu_{e})$.
Similarly, if there was any, 
CC MC templates made from an {\em oscillated} $\bar{\nu}_{e}$ event sample 
would be weighted with $P(\bar{\nu}_{\mu}\rightarrow \bar{\nu}_{e})$,
CC MC templates made from an {\em oscillated} ${\nu}_{\mu}$ event sample 
would be weighted with $P({\nu}_{e}\rightarrow {\nu}_{\mu})$, and
CC MC templates made from an {\em oscillated} $\bar{\nu}_{\mu}$ event sample 
would be weighted with $P(\bar{\nu}_{e}\rightarrow \bar{\nu}_{\mu})$.
Contributions from $\nu_{\tau}$ CC and $\bar{nu}_{\tau}$ CC are neglected as,
even when they are produced via oscillations, have a negligible effect on
observed distributions due to the high energy threshold for $\tau$ production.
The same oscillation parameters are used for both neutrino and anti-neutrino oscillations. 
Oscillations are not applied to the NC MC templates since, effectively, 
they serve as proxies for the mixtures of $\nu_{e} + \nu_{\mu} + \nu_{\tau}$ NC and
$\bar{\nu}_{e} + \bar{\nu}_{\mu} + \bar{\nu}_{\tau}$
MC templates which are unchanged under 3-flavour oscillations.


%
%
%

\subsubsection{Oscillations in a framework with N sterile neutrinos}
\label{sec:3pNosc}

%
%
%

\subsubsection{Non-standard interactions}
\label{sec:nsi}

As described in \ref{sec:3fosc}, neutrino oscillation probabilities are modified by forward coherent scattering of neutrinos off electrons in the matter they propagate through. Additional, flavour-dependent interactions could exist, leading to further modification of oscillation probabilities. These hypothetical interactions, which are usually referred to as Non-Standard Interactions (NSI), could modify the L/E dependence, or couple charged leptons and neutrinos of different flavours.

NSI are described by the following addition to the SM Lagrangian density
\begin{equation}
  \mathcal{L}_{NSI} = -2 \sqrt{2} G_{F} 
    \epsilon_{\alpha \beta}^{fc}
    \Big( \bar{f} \gamma^{\mu} P_{c} f \Big)
    \Big( \bar{\nu}_{\alpha} \gamma_{\mu} P_{L} \nu_{\beta} \Big)
\label{eq:nsi_lagrangian_density}
\end{equation}
where $\epsilon_{\alpha \beta}^{fc}$ are parameters expressing the size
of deviation from SM interactions for neutrinos of flavours $\alpha$,$\beta$,
and chirality $c$. Typically, in experimental NSI searches, the large parameter space
is reduced to a smaller set of effective parameters ($\epsilon_{\alpha \beta}$) by
summing $\epsilon_{\alpha \beta}^{fc}$ over chirality and fermion flavour
\begin{equation}
\epsilon_{\alpha \beta} = \sum_{f,c} \epsilon_{\alpha \beta}^{fc} \frac{N_f}{N_e}
\label{eq:nsi_epsilon_ab}
\end{equation}
where $N_f$ is the number density of fermion $f$, and
$N_e$ is the electron number density.

In the 3-flavour neutrino oscillation calculation, 
NSI effects enter as an additional contribution to the standard Hamiltonian shown in Eq.\ref{}. 
\begin{equation}
A_{matter} = 2EV \left(
\begin{array}{ccc}
1+\epsilon_{ee}  & \epsilon_{e\mu}^{\star} & \epsilon_{e\tau}^{\star} \\
\epsilon_{e\mu}  & \epsilon_{\mu\mu}       & \epsilon_{\mu\tau}^{\star} \\
\epsilon_{e\tau} & \epsilon_{\mu\tau}      & \epsilon_{\tau\tau}
\end{array} \right)
\end{equation}

%The new contribution has 6 new free parameters.

\subsection{Framework for parameterization of systematic effects}

Systematic effects enter via the term $R_{d;b;s;m}(r,t; \vec{f})$ in Eq. \ref{eq:valor_npred}.
Typically, but not always, the response $R_{d;b;s;m}(r,t; \vec{f})$ 
factorises and it can be written as 
\begin{equation}
  R_{d;b;s;m}(r,t; \vec{f}) = \prod_{i=0}^{N-1} R^{i}_{d;b;s;m}(r,t; f_i)
\label{eq:valor_response_factorization}
\end{equation}
For several systematics the response is linear and, therefore, 
\begin{equation}
  R^{i}_{d;b;s;m}(r,t; f_i) \propto f_i
\label{eq:valor_linear_response}
\end{equation}
For non-linear systematics, the response function $R^{i}_{d;b;s;m}(r,t; f_i)$ 
is usually pre-computed in the [-5$\sigma$, +5$\sigma$] range of the parameter $f_i$ 
and it is represented internally using Akima or cubic splines.
Values of systematic parameters that give a negative predicted number of events 
in any reconstructed bin in any interaction mode are not allowed.


\subsection{Statistical tools}
\label{S:valor_stat}

\subsubsection{Parameter estimation}

\subsubsection{Elimination of nuisance parameters}

\subsubsection{Construction of confidence intervals}

